\section{Conclusions}
\label{sec:conclusions}

We have built a search for heterotic line-bundle models in which every pruning
step is a theorem with stated assumptions, and used it to push the analysis from
the upstairs GUT stage through the equivariant descent to the holomorphic
superpotential.

On the side of \emph{search}, the central observation is that the slope
obstruction, used in earlier scans as a finite coefficient test, has an exact
and complete form---a fixed polyhedral cone meeting the span of the charge
vectors---that is monotone, prunes every branch, and is decided by a dual test
whose cost falls with depth. It removes a spurious candidate the finite test
keeps, and with it the search reproduces the published catalogs of both
benchmark geometries exactly. The slice decomposition that follows from the same
geometry turns 425 slices of the tetraquadric into theorems valid for
arbitrarily large charges, and isolates precisely the slices where only the
cubic growth of the index can bound the charges. Run over all 68 favourable
CICYs, the search never misses a published model, and an exact decision of
stability---validated by never rejecting one---shows that the published catalog
itself is incomplete: it lacks at least one poly-stable GUT model whose only
stability witness lies at a Kähler class with moduli ratio about 1:231.

On the side of \emph{physics}, computing the $\Gam$-representation on cohomology
from the group action lets the pipeline discover rather than verify: from our
own exhaustive survivor set it recovers all ten published standard-model orbits
on the tetraquadric quotient and finds eleven more. Their absence from the
published database is explained, in every case, by restrictions the published
scan states for itself. The superpotential machinery reproduces every operator
table of that database, exposes an internal inconsistency in almost half of its
down-type Yukawa tables, and---by evaluating leading-order couplings as genuine
triple products---detects vanishings that charge conservation cannot.

The answer to the phenomenological question is sobering and, we think, useful:
no model on this quotient, published or new, combines a leading-order top Yukawa
with rank-three Yukawa matrices and the absence of dangerous operators. The new
models fail in the same ways and at a comparable rate as the published ones. Two
of them (A04, A06) remain open, and deciding them requires exactly the
ingredients---higher-order couplings, singlet vacua, Kähler metrics---whose
absence we have documented rather than papered over.

The methods are not specific to the two benchmarks, as the run over all 68
favourable CICYs shows: the complete slope test needs only the triple
intersection numbers, the exact stability decision only the slopes, and the flow
formulation of the selection rules is geometry-independent. The natural next
steps are that survey at larger charges---where more models missing from the
catalog may appear---the exact treatment of rank-deficient charge matrices, the
extension of the triple-product evaluation to higher Koszul levels, and the
finiteness question of \cref{conj:C1}.
